\chapter{Agujeros Negros: Nodos de Compresión}
Los agujeros negros no son simples sumideros de materia, sino \emph{nodos de compresión} donde la información se almacena en el horizonte de eventos. En TCCU-6D, un agujero negro es una región donde el campo $\Phi$ alcanza un valor crítico, generando una superficie de coherencia perfecta.

\section{Métrica de Kerr-Schild}
Para modelar la curvatura extrema y el arrastre de marcos utilizamos la forma de Kerr-Schild:
\begin{equation}
g_{\mu\nu} = \eta_{\mu\nu} + 2H k_{\mu} k_{\nu}
\end{equation}
donde $H = \frac{GM}{c^2 r} \frac{r^3}{r^4 + a^2 z^2}$ para un agujero negro rotatorio con espín $a$.

\section{Termodinámica de los nodos de compresión}
La entropía de un agujero negro se reinterpreta como la cantidad de información comprimida. El principio holográfico emerge de manera natural: la información tridimensional se codifica en una superficie bidimensional. La coherencia del horizonte está dada por:
\begin{equation}
\Phi_{\text{H}} = \frac{A}{4G\hbar} \cdot \frac{1}{1 + e^{-\beta (T - T_0)}}
\end{equation}
Esta relación conecta la gravedad cuántica con la teoría de la información.
